\section{Conclusiones}
\label{sec:conclusiones}

Este trabajo implementó un simulador para redes descentralizadas y lo probó
con simulaciones de una Red de Expansión de Cobertura. Las conclusiones siguen
el orden de los objetivos específicos.

\textbf{Sobre los requerimientos (objetivo 1).} La revisión mostró que los
simuladores de propósito general evalúan una copia del protocolo, que los
emuladores ejecutan el programa real pero no traen un modelo listo de la pérdida
de señal entre pisos, y que las herramientas de redes entre pares no representan
el medio de radiofrecuencia. Además, cualquier simulador conoce toda la red, lo
que facilita decisiones que ningún nodo real podría tomar; el simulador de
partida lo hacía en tres lugares. De ahí salieron los diez requerimientos de la
Tabla~\ref{tab:requisitos}, entre ellos ejecutar el código del nodo, decidir
solo con información local y medir lo que cree cada nodo.

\textbf{Sobre el diseño (objetivo 2).} La arquitectura en capas deja el código
del nodo sin cambios y reemplaza únicamente la red por un modelo del medio que
depende de la distancia y de los pisos. Los paquetes avanzan con la tabla de
rutas de cada nodo y los Gateway se mueven solo con lo que oyen, con lo que se
eliminaron los tres usos de información global, que quedan solo como referencia
de comparación.

\textbf{Sobre la implementación (objetivo 3).} Se obtuvo un simulador que se
configura con archivos, con dos tipos de mapa, tres modos de movimiento de los
Gateway, tres de los usuarios, tres tipos de paquete (UDP, ICMP y TCP), caídas
programables y cuatro formas de ejecución, entre ellas una interfaz web en la
que cada función se puede observar. Su límite es el tamaño: corre más rápido
que el tiempo real hasta unas decenas de nodos, y más lento con 100 nodos
juntos. Su ventaja frente a la emulación es el control y las mediciones, no el
tamaño.

\textbf{Sobre las pruebas (objetivo 4).} De las once pruebas de simulación,
nueve dieron lo esperado y dos lo dieron en parte
(Tabla~\ref{tab:resumen-pruebas}). Con el despliegue de referencia y 10
semillas se concluye que:

\begin{itemize}
\item \textit{Cobertura.} La cobertura no muestra el daño a la conectividad:
la caída de un puente parte la malla 90~s con cobertura media de 1.000, y dos
fallos dirigidos reducen la entrega de datos de 0.502 a 0.284 sin una
diferencia de cobertura concluyente. Evaluar una red de este tipo exige medir
a la vez la física, la entrega de datos y lo que cree cada nodo.
\item \textit{Adaptabilidad.} La malla se puede adaptar sin coordinador: con
Gateway móviles, la partición por la caída del puente baja de 90~s a 6.2~s. El
costo es una cobertura de 0.922 en ese mismo fallo, rutas que tardan más en
rehacerse tras la caída de un nodo redundante y tablas de rutas menos exactas:
en el escenario \texttt{base}, la exactitud del conocimiento baja de 0.95 a
0.55.
\item \textit{Tipos de paquete.} Con el mismo tráfico, TCP entrega 0.924, UDP
0.772 e ICMP 0.646; TCP lo logra con cerca de 2.6 veces las transmisiones de
UDP.
\item \textit{Protocolos.} La comparación depende del lugar: la inundación
supera a B.A.T.M.A.N. en campo abierto con cerca del doble de transmisiones, y
queda muy por debajo en el edificio (0.430 frente a 0.786). Este resultado vale
para un medio sin colisiones.
\item \textit{Repetibilidad.} Las tablas de resultados salieron idénticas byte
a byte al repetir las corridas en dos computadores.
\end{itemize}

En conjunto, varios de estos resultados solo aparecieron al ejecutar el código
del nodo, exigir decisiones locales y medir lo que cree cada nodo, y algunos
contradijeron lo que el propio proyecto daba por cierto. Las conclusiones valen
para esta implementación bajo este medio. No se hizo una validación con equipos
físicos ni contra batman-adv, que es la principal tarea pendiente.

\section{Recomendaciones}
\label{sec:recomendaciones}

\begin{itemize}
\item \textbf{Validación externa.} Es la prioridad. Instalar el nodo de red
(\texttt{mesh/node.py}) en un grupo pequeño de equipos de radio, repetir las
pruebas de fallo y comparar con la simulación; y comparar una topología con
batman-adv en un emulador o con un modelo en ns-3.
\item \textbf{Acercar el protocolo a B.A.T.M.A.N.} Reenviar solo la copia del
mejor vecino y borrar los nodos de origen que llevan tiempo en silencio, y
medir cuánto se reducen los bucles y las pérdidas por rutas viejas.
\item \textbf{Mejorar el medio.} Añadir enlaces que no sean simétricos, cambios
en el tiempo, colisiones y una capa de acceso al medio, para responder a las
críticas conocidas a los modelos simples~\cite{Kotz2004Assumptions} y para
evaluar la inundación con su costo real~\cite{Ni1999BroadcastStorm}. Calibrar
el medio con mediciones en un edificio~\cite{Seidel1992Multifloor}.
\item \textbf{Retardo y transporte.} Resolver los saltos con un retardo en
lugar de dentro de un paso, lo que permitiría medir el retardo de los paquetes
y representar TCP con sus tiempos de espera y su control de
congestión~\cite{Eddy2022TCP}.
\item \textbf{Movimiento que tenga en cuenta la calidad.} Incluir en las reglas
la calidad de los enlaces y no solo el alcance, para que las tablas de rutas no
se queden atrás. Evaluar los umbrales en escenarios que no se usaron para
ajustarlos.
\item \textbf{Aprendizaje de las reglas locales.} El simulador ofrece un
entorno repetible, con agentes que solo tienen información parcial, resultados
medibles y un costo bajo por corrida. Esto permite plantear, como continuación
en inteligencia artificial, que los Gateway aprendan su regla de movimiento, y
compararla con la de este trabajo.
\item \textbf{Tamaño.} Optimizar el procesamiento de los OGM, que consume la
mayor parte del tiempo, para superar el límite actual de unos 100 nodos juntos.
\item \textbf{Estadística.} Aumentar el número de semillas en las comparaciones
de magnitud pequeña y aplicar una corrección por hacer varias
comparaciones~\cite{Jain1991Art}.
\end{itemize}
